\uuid{6DsG}
\exo7id{603}
\titre{exo7 603}
\auteur{gourio}
\organisation{exo7}
\datecreate{2001-09-01}
\isIndication{false}
\isCorrection{false}
\chapitre{Suite}
\sousChapitre{Suite récurrente linéaire}
\cpge{mpsi-mp2i.nombres-reels-suites.i, pcsi.nombres-reels-suites.g, ptsi.nombres-reels-suites.g, mpsi-mp2i.analyse-asymptotique.d, pcsi.analyse-asymptotique.d, mpsi-mp2i.analyse-asymptotique.b, pcsi.analyse-asymptotique.b, ptsi.analyse-asymptotique.b, bcpst1.suites-reelles}
\module{Analyse}
\niveau{L1}
\difficulte{}

\contenu{
\question{
\'Etudier la suite $(u_{n})_{n\in \Nn}$ d\'{e}finie par :
$$u_{0}\in ]0,\frac{\pi }{2}[,u_{n+1}=\sin u_{n}.$$
Donner $\lim_{x\rightarrow 0}\frac{1}{\sin ^{2}x}-\frac{1}{x^{2}}$
,(r\'{e}ponse: $\frac{1}{3}$) en d\'{e}duire un \'{e}quivalent de $u_{n}^{-2} $
donc de $u_{n}.$
}
}
